% !TeX root = ../main.tex

\section{Referencial teórico}

\subsection{Cadastro Ambiental Rural e Sistema de Cadastro Ambiental Rural}

A Lei n. 12.651/2012 instituiu o CAR no âmbito do Sistema Nacional de Informação sobre Meio Ambiente. O cadastro reúne informações do perímetro do imóvel rural, das Áreas de Preservação Permanente, das áreas de uso restrito, das áreas consolidadas, das Reservas Legais e dos remanescentes de vegetação nativa \cite{brasil2012}. A organização dessas informações em uma base nacional amplia as possibilidades de controle e planejamento ambiental, mas a natureza declaratória da inscrição exige análise posterior pelo poder público.

O Decreto n. 7.830/2012 regulamentou o SICAR como sistema eletrônico nacional destinado ao gerenciamento das informações ambientais dos imóveis rurais. Entre seus objetivos estão receber e integrar os dados do CAR, controlar e acompanhar a manutenção e a regularização da vegetação nativa e disponibilizar informações públicas de natureza ambiental \cite{brasil2012decreto}. CAR e SICAR, portanto, não são sinônimos: o primeiro é o registro ambiental do imóvel, enquanto o segundo corresponde à infraestrutura de informação que recebe, organiza e permite analisar as declarações.

A emissão do recibo comprova que a inscrição foi transmitida ao sistema, mas não representa aprovação da localização da Reserva Legal nem reconhecimento de propriedade ou posse. A análise e a validação cabem ao órgão ambiental estadual competente \cite{car2026}. Essa distinção é essencial para o presente estudo, pois o cruzamento cartográfico identifica coerência ou divergência entre bases, sem produzir decisão sobre regularidade ambiental.

As informações geoespaciais declaradas são produzidas em condições heterogêneas. A delimitação pode resultar de levantamento de campo, interpretação de imagens ou vetorização sobre bases com diferentes escalas e datas. Além disso, retificações posteriores podem alterar as geometrias disponíveis no sistema. Essas características, associadas à natureza autodeclaratória do cadastro, exigem o registro da data de aquisição, a verificação do sistema de referência e o controle de sobreposições, duplicidades e inconsistências geométricas antes do cálculo de áreas \cite{car2026,matias2024}.

\subsection{Estrutura das camadas geoespaciais do CAR}

A inscrição organiza dados documentais e geoespaciais em conjuntos temáticos. A camada de área do imóvel representa o perímetro da propriedade ou posse rural e funciona como limite de referência para as demais feições. A geometria declarada no CAR tem finalidade ambiental, não constitui título para o reconhecimento do direito de propriedade ou posse e não substitui os procedimentos fundiários de cadastramento e certificação do imóvel \cite[art. 29, \S~2º]{brasil2012}.

Os remanescentes de vegetação nativa correspondem aos fragmentos declarados com cobertura vegetal natural. Essa camada pode se sobrepor a APP e Reserva Legal, pois os temas descrevem atributos diferentes de uma mesma área: cobertura existente, proteção decorrente da localização e destinação jurídica. A Instrução Normativa MMA n. 2/2014 determina que remanescentes existentes após 22 de julho de 2008 não perdem essa classificação em razão de intervenção não autorizada \cite{mma2014}.

A utilização das informações do CAR em análises territoriais exige atenção à qualidade das geometrias declaradas. As delimitações podem apresentar sobreposições, duplicidades e inconsistências topológicas que interferem diretamente nos cálculos de área e nas operações de geoprocessamento. \citeonline{matias2024} discutem essas limitações e apresentam procedimentos voltados ao tratamento de sobreposições em imóveis rurais cadastrados.

As dificuldades de integração de bases territoriais também são discutidas por \citeonline{sparovek2019}, que analisaram informações fundiárias brasileiras e identificaram sobreposições entre diferentes categorias de registro territorial. Embora a pesquisa apresente uma abordagem fundiária mais ampla, seus resultados evidenciam a importância da consistência espacial na utilização de bases cadastrais.

Nesse sentido, a análise de remanescentes de vegetação nativa declarados deve considerar a possibilidade de coincidência geométrica entre diferentes feições. A união dos polígonos permite quantificar a superfície espacialmente distinta do conjunto analisado, sem somar repetidamente as parcelas coincidentes, conforme o problema topológico discutido por \citeonline{matias2024}.

As principais camadas ambientais declaradas no CAR, bem como suas respectivas características e finalidades, estão apresentadas na Tabela~\ref{tab:camadascar}.

\begin{table}[H]
\centering
\caption{Principais camadas geoespaciais declaradas no CAR}
\label{tab:camadascar}

\begingroup
\small
\SingleSpacing

% Espaçamento vertical entre as linhas
\renewcommand{\arraystretch}{1.35}

% Espaçamento horizontal entre as colunas
\setlength{\tabcolsep}{8pt}

\begin{tabularx}{\textwidth}{
>{\RaggedRight\arraybackslash}p{4.1cm}
>{\arraybackslash}X
}
\toprule
\textbf{Camada} & \textbf{Conteúdo representado} \\
\midrule

Área do imóvel rural &
Perímetro da propriedade ou posse usado como referência espacial do cadastro \\

Remanescente de vegetação nativa &
Fragmentos declarados de vegetação nativa existentes no imóvel. \\

Área de Preservação Permanente &
Áreas protegidas em razão de sua localização, como margens de cursos d'água, entorno de nascentes e outras hipóteses legais. \\

Reserva Legal &
Área localizada no interior do imóvel destinada à conservação e ao uso sustentável, declarada como existente, proposta ou vinculada a situação administrativa específica. \\

Área de uso restrito &
Espaços submetidos a regras próprias de uso e conservação previstas na legislação. \\

Área rural consolidada &
Área com ocupação antrópica preexistente a 22 de julho de 2008, nos termos da Lei n. 12.651/2012. \\

Servidão administrativa e utilidade pública &
Faixas ou áreas afetadas por infraestrutura e outras finalidades públicas que interferem na área disponível do imóvel. \\

Hidrografia e nascentes &
Feições lineares, poligonais ou pontuais usadas na caracterização ambiental e na delimitação de áreas protegidas. \\

\bottomrule
\end{tabularx}

\vspace{0.25cm}
\fontefigura{Elaboração própria (2026) com informações de Brasil (2012, 2014) e SICAR.}

\endgroup
\end{table}

As camadas ambientais do CAR não representam necessariamente categorias espaciais independentes, uma vez que uma mesma área pode estar associada a diferentes classificações. Por exemplo, um fragmento de vegetação nativa pode estar localizado em Área de Preservação Permanente (APP) e, simultaneamente, integrar a Reserva Legal (RL) do imóvel. Dessa forma, a interpretação e a quantificação dessas camadas devem considerar a possibilidade de interseções decorrentes de funções ambientais e jurídicas distintas, evitando que a coincidência espacial seja tratada automaticamente como inconsistência cadastral.

\subsection{Etapas de inscrição, análise e retificação}

O cadastramento inicia-se com a identificação do cadastrante e do proprietário ou possuidor, seguida da inclusão dos documentos de comprovação da propriedade ou posse. Na etapa geoespacial são informados o perímetro do imóvel e, quando existentes, os remanescentes de vegetação nativa, as APP, as áreas de uso restrito, as áreas consolidadas e a Reserva Legal \cite{mma2014,car2026}. O módulo pré-preenchido pode apresentar informações recuperadas de bases oficiais, mas o declarante deve revisá-las, complementá-las e confirmá-las.

Depois da conferência, os dados são gravados e enviados ao SICAR. A transmissão gera o Recibo de Inscrição, documento que comprova a entrega das informações. O acompanhamento das etapas de inscrição, análise e regularização ocorre por meio do Demonstrativo da Situação do CAR e da Central do Proprietário ou Possuidor \cite{brasilservicocar2025}. A situação administrativa do cadastro não pode ser deduzida apenas a partir das camadas vetoriais públicas.

Na análise, o órgão estadual confronta a declaração com bases de referência e verifica a consistência das informações. Havendo divergências ou documentação insuficiente, o responsável pode ser notificado para prestar esclarecimentos ou retificar o cadastro. Quando forem identificados passivos ambientais, o processo pode prosseguir para os instrumentos de regularização previstos na legislação, inclusive o Programa de Regularização Ambiental \cite{sfbregularizacao2026}. A retificação altera os dados documentais ou geoespaciais e deve ser considerada na definição da versão da base utilizada em pesquisas.

Esse fluxo explica por que duas extrações do SICAR realizadas em datas diferentes podem apresentar quantidades e geometrias distintas. A rastreabilidade depende de conservar a data do download e os arquivos efetivamente analisados. Também explica por que a divergência com o MapBiomas deve ser tratada como indicação para exame, e não como confirmação automática de erro na declaração.

\subsection{Operacionalização do CAR na Paraíba}

Na Paraíba, a Superintendência de Administração do Meio Ambiente (SUDEMA) atua como órgão ambiental estadual responsável pelos procedimentos relacionados ao CAR. A página institucional reúne formulários para requerimento de análise, alteração de titularidade e endereço eletrônico, cancelamento do cadastro, regularização de Reserva Legal e atendimento por responsável técnico. Também disponibiliza os dispositivos normativos estaduais e fluxos de regularização ambiental \cite{sudema2026}.

A SUDEMA esclarece que a inscrição constitui o primeiro passo da regularização ambiental e atribui ao proprietário ou possuidor a responsabilidade pelas informações declaradas e por suas retificações. Quando a análise identifica pendências ou inconsistências, o responsável deve fornecer informações complementares ou promover as correções solicitadas \cite{sudema2026}. O processo administrativo estadual, portanto, envolve comunicação com o declarante e avaliação de outras informações além da geometria publicada.

A Paraíba utiliza a solução tecnológica do SICAR disponibilizada pelo Serviço Florestal Brasileiro, enquanto o atendimento relativo à inscrição, à análise e ao cancelamento deve ser realizado pelo órgão estadual competente \cite{brasilservicocar2025}. A base vetorial pública representa uma extração do sistema na data de aquisição e não reúne, necessariamente, todos os documentos e registros administrativos utilizados pela SUDEMA na análise individual.

Essa divisão de competências fundamenta o limite interpretativo do estudo. O Serviço Florestal Brasileiro apoia a gestão e a integração nacional do SICAR, ao passo que a avaliação administrativa ocorre no âmbito estadual. Assim, os resultados espaciais podem subsidiar a triagem e a priorização de verificações, mas não substituem a análise conduzida pela SUDEMA nem os procedimentos de regularização ambiental aplicáveis a cada imóvel.


\subsection{Sensoriamento remoto e mapeamento da cobertura e uso da terra}

O sensoriamento remoto permite caracterizar a cobertura da superfície terrestre a partir de dados registrados por sensores orbitais. Associados aos Sistemas de Informação Geográfica (SIG), esses dados podem ser classificados e analisados para representar a distribuição espacial da vegetação, das atividades agropecuárias e das áreas urbanizadas. A interpretação dos produtos resultantes depende, entretanto, das características das imagens, da legenda temática e dos procedimentos de classificação e validação \cite{foody2002,souza2020}.

Em estudos de bacias hidrográficas, a análise multitemporal da cobertura da terra permite examinar mudanças na ocupação territorial e sua relação com processos ambientais. A comparabilidade entre períodos, contudo, depende da compatibilidade entre sensores, resoluções espaciais, datas de aquisição, legendas e métodos de classificação \cite{foody2002,soares2024}.

Na bacia hidrográfica do rio Cuiá, em João Pessoa, \citeonline{soares2024} utilizou imagens Landsat e PlanetScope para caracterizar a cobertura e o uso da terra em diferentes períodos. A classificação supervisionada, realizada pelo método da máxima verossimilhança, resultou em oito classes temáticas, posteriormente empregadas na análise dos efeitos da urbanização sobre o escoamento superficial. Embora tenha finalidade hidrológica distinta da presente pesquisa, o trabalho evidencia a aplicação de dados orbitais na caracterização ambiental de bacias costeiras paraibanas.

No Brasil, o projeto MapBiomas disponibiliza séries anuais de cobertura e uso da terra, produzidas a partir de imagens orbitais e procedimentos de classificação. \citeonline{souza2020} descrevem a utilização do acervo Landsat e do processamento em nuvem para reconstruir mudanças na cobertura da terra dos biomas brasileiros ao longo de aproximadamente três décadas.

A resolução espacial condiciona a representação das feições mapeadas. Em produtos com pixels de 30 m, fragmentos estreitos e pequenas áreas de vegetação podem não ser individualizados, sobretudo quando diferentes coberturas ocorrem em uma mesma célula. A série MapBiomas utiliza imagens Landsat e resolução espacial nominal de 30 m, de modo que a interpretação de seus resultados deve considerar a unidade espacial do produto e a escala da análise \cite{souza2020,mapbiomas2026}.

A utilização de diferentes fontes de informação geográfica também exige atenção à compatibilidade entre suas representações. No caso do CAR e do MapBiomas, a comparação envolve polígonos delimitados pelos declarantes e classes de cobertura da terra provenientes de classificação de imagens orbitais, que podem apresentar diferenças espaciais e temáticas.



\subsection{Comparação entre dados declaratórios e mapeamento por sensoriamento remoto}

O Cadastro Ambiental Rural e o MapBiomas representam informações ambientais produzidas por procedimentos distintos. Enquanto o CAR reúne feições geoespaciais autodeclaradas, o MapBiomas apresenta classes de cobertura e uso da terra obtidas por classificação de imagens orbitais. A comparação entre essas bases permite identificar áreas de correspondência e divergência espacial, mas não estabelece, isoladamente, qual representação corresponde às condições reais do território.

A compatibilidade entre mapas depende da definição das classes utilizadas e dos procedimentos de classificação. \citeonline{neves2020}, ao compararem os produtos TerraClass e MapBiomas na Amazônia brasileira, demonstraram que diferenças entre legendas influenciam os resultados de concordância temática e espacial.

\citeonline{pontius2011} discutem a importância de distinguir divergências relacionadas à quantidade de área atribuída às classes daquelas associadas à sua distribuição espacial. Essa abordagem demonstra que mapas podem apresentar proporções semelhantes de determinadas coberturas, mas diferir quanto à localização das feições representadas.

Outro aspecto importante é a distinção entre concordância espacial e acurácia temática. \citeonline{foody2002} discute as limitações dos procedimentos de avaliação de classificações de cobertura da terra, enquanto \citeonline{olofsson2014} apresentam recomendações para estimativas de área e avaliação de acurácia com dados independentes de referência.

\citeonline{olofsson2013} também destacam que os erros de classificação podem influenciar as estimativas de área obtidas diretamente de mapas temáticos. Dessa forma, a quantificação da correspondência entre duas bases não equivale necessariamente à validação de suas classificações.

À luz dessas abordagens, a sobreposição entre o CAR e o MapBiomas expressa convergência espacial, e não acurácia. A avaliação formal de acurácia exige amostragem probabilística e dados de referência independentes \cite{olofsson2014}. Essa convergência tampouco substitui a análise administrativa pelo órgão estadual competente \cite{car2026}.


\newpage
